\documentclass[12pt,a4paper]{article}
\usepackage[utf8]{inputenc}
\usepackage{amsmath, amssymb, amsthm}
\usepackage[ngerman]{babel}
\usepackage[T1]{fontenc}
\usepackage{hyperref}
\usepackage[ngerman]{cleveref}
\usepackage{xurl}
\usepackage{parskip}
\hypersetup{hidelinks}

% Include Pictures
\usepackage{graphicx, float}
\graphicspath{{images/}}

% Page setup
\usepackage{geometry}
\geometry{
left=3cm,
right=3cm,
top=3cm,
bottom=3cm
}


% Line hight
\renewcommand{\baselinestretch}{1.25}

% Parskip and parindent
\setlength{\parindent}{0pt}
\setlength{\parskip}{1ex}


\begin{document}

% ------------------------------------------------
% Titelseite
% ------------------------------------------------
\begin{titlepage}
\centering

% Kopf
\vspace*{1.0cm}

{\large Universität Freiburg\par}
{\large Physiklabor für Anfänger*innen 1\par}
{\large Ferienpraktikum im Sommersemester 2026\par}

\vspace{2.5cm}

% Versuchstitel
{\Huge\bfseries Versuch 18\\ Reversionspendel\par}

\vspace{3.5cm}

% Namen
{\large Lorenz Klug, Matti Jacob, Oskar Gentner\par}
{\large (Gruppe 143)\par}

\vspace{2.0cm}

% Datum
{\large 18.09.2026\par}

\vspace{1.5cm}

{\large Datum der Durchführung: 17.09.2026\par}

\vspace{1.5cm}

{\large Assistent*in: Johanna Böck\par}

\vfill

\end{titlepage}


% ------------------------------------------------
% Inhaltsverzeichnis
% ------------------------------------------------
\tableofcontents
\newpage

% Beginn des Dokuments


\section{Ziel des Versuchs}

\section{Versuchsdurchführung}
  \begin{figure}[H]
    \centering
    \includegraphics[width=\textwidth]{images/}
    \caption{caption}
    \label{fig:}
  \end{figure}

\section{Auswertung und Fehleranalyse}

\subsection{Umgebungsbedingungen}
 
\subsection{Fehlerfortpflanzung}

Wir verwenden die Gauß'sche Fehlerfortpflanzung um den Fehler von $g$ zu bestimmen.

\section{Diskussion}
  
  \subsection{Ergebnisse}


  \subsection{Vergleich mit dem Literaturwert}

Wir führen den z-Test auf Verträglichkeit unseres Ergebnisses $g_B$ mit dem Literaturwert $g_L$ durch.
  
\subsection{Störeinflüsse}

  \begin{itemize}
    \item 
 \end{itemize}
    
  \subsection{Verbesserungsvorschläge}
  Wir haben einige Verbesserungsvorschläge um eine präzisere Messung zu ermöglichen gesammelt:
  \begin{itemize}
    \item 
  \end{itemize}
    
\section{Autorenschaft}
  Alle Autoren haben zu allen Inhalten des Protokolls zu gleichen Teilen beigetragen.
    
\section{Anhang}

\newpage
\section{Labornotizen}


\end{document}
